%
% 24-liouville.tex -- Der Satz von Liouville
%
% (c) 2026 Prof Dr Andreas Müller
%
\subsection{Der Satz von Liouville}
Die Cauchy-Formel für die Ableitungen hat die unerwartete Konsquenz,
dass auf ganz $\mathbb{C}$ definierte holomorphe Funktionen, die nicht
Polynome sind, exponentiell schnell anwachsen müssen, wenn $z$ gross wird.

\begin{definition}[polynomiell beschränkt]
Eine in ganz $\mathbb{C}$ definierte komplexe Funktion $f(z)$ heisst
\emph{polynomiell beschränkt}, wenn es $N\in\mathbb{N}$ und $a>0$
\index{polynomiell beschränkt}%
gibt derart, dass
\begin{equation}
|f(z)|
\le a(1+|z|^N)
\label{buch:integration:eqn:liouvilleschranke}
\end{equation}
für alle $z\in\mathbb{C}$ gilt.
\end{definition}

\begin{satz}[Liouville]
\label{buch:integration:satz:liouville}
\index{Satz!von Liouville}%
\index{Liouville, Joseph}%
Sei $f$ eine auf ganz $\mathbb{C}$ holomorphe Funktion, die polynomiell
beschränkt ist.
Dann ist $f$ ein Polynom vom Grad höchstens $N$.
\end{satz}

\begin{proof}
Wir wenden den Satz
\ref{buch:integration:satz:ableitungn}
für Ableitungen einer holomorphen Funktion auf den Weg
$\gamma(t)=z_0+Re^{2\pi it}$.
Nach der Formel \eqref{buch:integration:eqn:ableitungn}
gilt für für alle $k$
\begin{align*}
f^{(k)}(z_0)
&=
\frac{k!}{2\pi i}
\oint_\gamma
\frac{f(z)}{(z-z_0)^{k+1}}
\,dz.
\intertext{Mit der Voraussetzung
\eqref{buch:integration:eqn:liouvilleschranke}
wir dies durch}
|f^{(k)}(z_0)|
&=
\frac{k!}{2\pi}
\oint_\gamma
\biggl|
\frac{f(z)}{(z-z_0)^{k+1}}
\biggr|
\,dz
\\
&\le
\frac{k!}{2\pi}
\int_0^1
\frac{|f(\gamma(t))|}{R^{k+1}}
2\pi R
\,dt
\\
&\le
k!
\int_0^1
\frac{|f(\gamma(t))|}{R^k}
\,dt
\\
&\le
\frac{k!}{R^k}
\int_0^1
a(1+|\gamma(t)|^N)
\,dt
\\
&\le
\frac{k!}{R^k}
\int_0^1
a(1+(|z_0|+R)^N)
\,dt
\\
&\le
k!
a
\frac{
1+(|z_0|+R)^N
}{R^k}
\end{align*}
abgeschätzt.
Für $R\to\infty$ strebt der Bruch auf der rechten Seite für $k>N$ 
gegen $0$, so dass $f^{(k)}(z_0)=0$ folgt.
Alle Ableitungen höherer als $N$-ter Ordnung der Funktion $f$ verschwinden,
also ist $f$ ein Polynom vom Grad höchstens $N$.
\end{proof}

Man beachte, dass der Grenzübergang im letzten Schritt des Beweises 
verwendet, dass die Funktion auf ganz $\mathbb{C}$ holomorph ist.
Sobald es Punkte in $\mathbb{C}$ gibt, in denen die Funktion nicht
holomorph ist, gilt auch die Integralformel nicht mehr und damit
bricht die Abschätzungskette des Beweises zusammen.

Der Spezialfall $N=0$ liefert die folgende Aussage.

\begin{korollar}
Eine beschränkte, auf ganz $\mathbb{C}$ holomorphe Funktion ist
konstant.
\end{korollar}

Für reelle Funktionen kann ein solches Resultat nicht gelten.
Die trigonometrischen Funktionen $\sin x$ und $\cos x$ sind
für alle reellen Argumente von $x$ beschränkt.
In $\mathbb{R}$ gilt für beide Funktionen
\[
|f(x)| \le 1,
\]
was der Ungleichung
\eqref{buch:integration:eqn:liouvilleschranke}
mit $N=0$ und $a=1$ entspricht.
Die trigonometrischen Funktionen sind aber alles andere als konstant
und auch keine Polynome, da Ableitungen beliebig hoher Ordnung nicht
verschwinden.

Umgekehrt bedeutet der Satz von Liouville, dass eine Funktion,
die durch eine für alle $z\in\mathbb{C}$ konvergente Potenzreihe
definiert ist, die kein Polynom ist, nicht polynomiell beschränkt
sein kann.
Für die Sinus- und Kosinusfunktionen zeigt die Definition
\[
\cos z = \frac{e^{iz}+e^{-iz}}{2}
\qquad\text{und}\qquad
\sin z = \frac{e^{z}-e^{-iz}}{2i},
\]
dass sich die Funktionen für imaginäre Argumente wie Exponentialfunktionen
verhalten, die nicht polynomiell beschränkt sind.

Holomorphe Funktionen, die keine Polynome sind, aber für grosse $z$
nur polynomiell anwachsen, können nicht in ganz $\mathbb{C}$
holomorph sein.
Rationale Funktionen, bei denen der Nennergrad grösser ist als der
Zählergrad, sind Beispiele für dieses Phänomen.
Tatsächlich hat eine solche Funktion immer mindestens eine Polstelle,
in der die Funktion nicht definiert ist.
